\documentclass{article}
\usepackage{amsmath, amssymb, ctex, graphicx, color}
\usepackage[margin=1in]{geometry}
\title{第5章-单层网络-分类 -- 第5.3节-生成式分类器}
\author{《深度学习基础》课程小组}
% \date{}
 
\begin{document}
\maketitle
\tableofcontents

\section{Generative Classifiers}

We turn next to a probabilistic view of classification and show how models with linear decision boundaries arise from simple assumptions about the distribution of the data. We have already discussed the distinction between the discriminative and the generative approaches to classification. Here we will adopt a generative approach in which we model the class-conditional densities $p(\mathbf{x}|\mathcal{C}_{k})$ as well as the class priors $p(\mathcal{C}_{k})$ and then use these to compute posterior probabilities $p(\mathcal{C}_{k}|\mathbf{x})$ through Bayes' theorem.

First, consider problems having two classes. The posterior probability for class $\mathcal{C}_{1}$ can be written as

\begin{equation}
p(\mathcal{C}_{1}|\mathbf{x}) = \frac{p(\mathbf{x}|\mathcal{C}_{1})p(\mathcal{C}_{1})}{p(\mathbf{x}|\mathcal{C}_{1})p(\mathcal{C}_{1})+p(\mathbf{x}|\mathcal{C}_{2})p(\mathcal{C}_{2})} = \frac{1}{1+\exp(-a)}=\sigma(a)
\tag{5.40}
\end{equation}

where we have defined

\begin{equation}
a=\ln\frac{p(\mathbf{x}|\mathcal{C}_{1})p(\mathcal{C}_{1})}{p(\mathbf{x}|\mathcal{C}_{2})p(\mathcal{C}_{2})}
\tag{5.41}
\end{equation}

and $\sigma(a)$ is the logistic sigmoid function defined by

\begin{equation}
\sigma(a)=\frac{1}{1+\exp(-a)},
\tag{5.42}
\end{equation}

which is plotted in Figure 5.12. The term 'sigmoid' means S-shaped. This type of function is sometimes also called a 'squashing function' because it maps the whole real axis into a finite interval. The logistic sigmoid has been encountered already in earlier chapters and plays an important role in many classification algorithms. It satisfies the following symmetry property:

\begin{equation}
\sigma(-a)=1-\sigma(a)
\tag{5.43}
\end{equation}

as is easily verified. The inverse of the logistic sigmoid is given by

\begin{equation}
a=\ln\left(\frac{\sigma}{1-\sigma}\right)
\tag{5.44}
\end{equation}

and is known as the logit function. It represents the log of the ratio of probabilities 
$$
\ln[p(\mathcal{C}_{1}|\mathbf{x})/p(\mathcal{C}_{2}|\mathbf{x})]
$$ 
for the two classes, also known as the log odds.

Note that in (5.40), we have simply rewritten the posterior probabilities in an equivalent form, and so the appearance of the logistic sigmoid may seem artificial. However, it will have significance provided $a(\mathbf{x})$ has a constrained functional form. We will shortly consider situations in which $a(\mathbf{x})$ is a linear function of $\mathbf{x}$, in which case the posterior probability is governed by a generalized linear model.

If there are $K>2$ classes, we have

\begin{equation}
p(\mathcal{C}_{k}|\mathbf{x}) = \frac{p(\mathbf{x}|\mathcal{C}_{k})p(\mathcal{C}_{k})}{\sum_{j}p(\mathbf{x}|\mathcal{C}_{j})p(\mathcal{C}_{j})} = \frac{\exp(a_{k})}{\sum_{j}\exp(a_{j})},
\tag{5.45}
\end{equation}

which is known as the normalized exponential and can be regarded as a multi-class generalization of the logistic sigmoid. Here the quantities $a_{k}$ are defined by

\begin{equation}
a_{k}=\ln\left(p(\mathbf{x}|\mathcal{C}_{k})p(\mathcal{C}_{k})\right).
\tag{5.46}
\end{equation}

The normalized exponential is also known as the softmax function, as it represents a smoothed version of the 'max' function because, if $a_{k}\gg a_{j}$ for all $j\neq k$, then $p(\mathcal{C}_{k}|\mathbf{x})\simeq 1$, and $p(\mathcal{C}_{j}|\mathbf{x})\simeq 0$.

We now investigate the consequences of choosing specific forms for the class-conditional densities, looking first at continuous input variables $\mathbf{x}$ and then discussing briefly discrete inputs.

\subsection{Continuous Inputs}

Let us assume that the class-conditional densities are Gaussian. We will then explore the resulting form for the posterior probabilities. To start with, we will assume that all classes share the same covariance matrix $\boldsymbol{\Sigma}$. Thus, the density for class $\mathcal{C}_{k}$ is given by

\begin{equation}
p(\mathbf{x}|\mathcal{C}_{k})=\frac{1}{(2\pi)^{D/2}}\frac{1}{|\boldsymbol{\Sigma}|^{1/2}}\exp\left\{-\frac{1}{2}(\mathbf{x}-\boldsymbol{\mu}_{k})^{\mathrm{T}}\boldsymbol{\Sigma}^{-1}(\mathbf{x}-\boldsymbol{\mu}_{k})\right\}.
\tag{5.47}
\end{equation}

First, suppose that we have two classes. From (5.40) and (5.41), we have

\begin{equation}
p(\mathcal{C}_{1}|\mathbf{x})=\sigma(\mathbf{w}^{\mathrm{T}}\mathbf{x}+w_{0})
\tag{5.48}
\end{equation}

where we have defined

\begin{align}
\mathbf{w} &= \boldsymbol{\Sigma}^{-1}(\boldsymbol{\mu}_{1}-\boldsymbol{\mu}_{2}) \tag{5.49}\\ 
w_{0} &= -\frac{1}{2}\boldsymbol{\mu}_{1}^{\mathrm{T}}\boldsymbol{\Sigma}^{-1}\boldsymbol{\mu}_{1}+\frac{1}{2}\boldsymbol{\mu}_{2}^{\mathrm{T}}\boldsymbol{\Sigma}^{-1}\boldsymbol{\mu}_{2}+\ln\frac{p(\mathcal{C}_{1})}{p(\mathcal{C}_{2})}. \tag{5.50}
\end{align}

We see that the quadratic terms in $\mathbf{x}$ from the exponents of the Gaussian densities have cancelled (due to the assumption of common covariance matrices), leading to a linear function of $\mathbf{x}$ in the argument of the logistic sigmoid. This result is illustrated for a two-dimensional input space $\mathbf{x}$ in Figure 5.13. The resulting decision boundaries correspond to surfaces along which the posterior probabilities $p(\mathcal{C}_{k}|\mathbf{x})$ are constant and so will be given by linear functions of $\mathbf{x}$, and therefore the decision boundaries are linear in input space. The prior probabilities $p(\mathcal{C}_{k})$ enter only through the bias parameter $w_{0}$, so that changes in the priors have the effect of making parallel shifts of the decision boundary and more generally of the parallel contours of constant posterior probability.

For the general case of $K$ classes, the posterior probabilities are given by (5.45) where, from (5.46) and (5.47), we have

\begin{equation}
a_{k}(\mathbf{x})=\mathbf{w}_{k}^{\mathrm{T}}\mathbf{x}+w_{k0}
\tag{5.51}
\end{equation}

in which we have defined

\begin{align}
\mathbf{w}_{k} &= \boldsymbol{\Sigma}^{-1}\boldsymbol{\mu}_{k} \tag{5.52}\\ 
w_{k0} &= -\frac{1}{2}\boldsymbol{\mu}_{k}^{\mathrm{T}}\boldsymbol{\Sigma}^{-1}\boldsymbol{\mu}_{k}+\ln p(\mathcal{C}_{k}). \tag{5.53}
\end{align}

We see that the $a_{k}(\mathbf{x})$ are again linear functions of $\mathbf{x}$ as a consequence of the cancellation of the quadratic terms due to the shared covariances. The resulting decision boundaries, corresponding to the minimum misclassification rate, will occur when two of the posterior probabilities (the two largest) are equal, and so will be defined by linear functions of $\mathbf{x}$. Thus, we again have a generalized linear model.

If we relax the assumption of a shared covariance matrix and allow each class-conditional density $p(\mathbf{x}|\mathcal{C}_{k})$ to have its own covariance matrix $\boldsymbol{\Sigma}_{k}$, then the earlier cancellations will no longer occur, and we will obtain quadratic functions of $\mathbf{x}$, giving rise to a quadratic discriminant. The linear and quadratic decision boundaries are illustrated in Figure 5.14.

\subsection{Maximum Likelihood Solution}

Once we have specified a parametric functional form for the class-conditional densities $p(\mathbf{x}|\mathcal{C}_{k})$, we can then determine the values of the parameters, together with the prior class probabilities $p(\mathcal{C}_{k})$, using maximum likelihood. This requires a data set comprising observations of $\mathbf{x}$ along with their corresponding class labels.

First, suppose we have two classes, each having a Gaussian class-conditional density with a shared covariance matrix, and suppose we have a data set $\{\mathbf{x}_{n},t_{n}\}$ where $n=1,\ldots,N$. Here $t_{n}=1$ denotes class $\mathcal{C}_{1}$ and $t_{n}=0$ denotes class $\mathcal{C}_{2}$. We denote the prior class probability $p(\mathcal{C}_{1})=\pi$, so that $p(\mathcal{C}_{2})=1-\pi$. For a data point $\mathbf{x}_{n}$ from class $\mathcal{C}_{1}$, we have $t_{n}=1$ and hence

\begin{equation}
p(\mathbf{x}_{n},\mathcal{C}_{1})=p(\mathcal{C}_{1})p(\mathbf{x}_{n}|\mathcal{C}_{1})=\pi\mathcal{N}(\mathbf{x}_{n}|\boldsymbol{\mu}_{1},\boldsymbol{\Sigma}).
\end{equation}

Similarly for class $\mathcal{C}_{2}$, we have $t_{n}=0$ and hence

\begin{equation}
p(\mathbf{x}_{n},\mathcal{C}_{2})=p(\mathcal{C}_{2})p(\mathbf{x}_{n}|\mathcal{C}_{2})=(1-\pi)\mathcal{N}(\mathbf{x}_{n}|\boldsymbol{\mu}_{2},\boldsymbol{\Sigma}).
\end{equation}

Thus, the likelihood function is given by

\begin{equation}
p(\mathbf{t},\mathbf{X}|\pi,\boldsymbol{\mu}_{1},\boldsymbol{\mu}_{2},\boldsymbol{\Sigma})=\prod_{n=1}^{N}\left[\pi\mathcal{N}(\mathbf{x}_{n}|\boldsymbol{\mu}_{1},\boldsymbol{\Sigma})\right]^{t_{n}}\left[(1-\pi)\mathcal{N}(\mathbf{x}_{n}|\boldsymbol{\mu}_{2},\boldsymbol{\Sigma})\right]^{1-t_{n}}
\tag{5.54}
\end{equation}

where $\mathbf{t}=(t_{1},\ldots,t_{N})^{\mathrm{T}}$. As usual, it is convenient to maximize the log of the likelihood function. Consider first the maximization with respect to $\pi$. The terms in the log likelihood function that depend on $\pi$ are

\begin{equation}
\sum_{n=1}^{N}\left\{t_{n}\ln\pi+(1-t_{n})\ln(1-\pi)\right\}.
\tag{5.55}
\end{equation}

Setting the derivative with respect to $\pi$ equal to zero and rearranging, we obtain

\begin{equation}
\pi=\frac{1}{N}\sum_{n=1}^{N}t_{n}=\frac{N_{1}}{N}=\frac{N_{1}}{N_{1}+N_{2}}
\tag{5.56}
\end{equation}

where $N_{1}$ denotes the total number of data points in class $\mathcal{C}_{1}$, and $N_{2}$ denotes the total number of data points in class $\mathcal{C}_{2}$. Thus, the maximum likelihood estimate for $\pi$ is simply the fraction of points in class $\mathcal{C}_{1}$ as expected. This result is easily generalized to the multi-class case where again the maximum likelihood estimate of the prior probability associated with class $\mathcal{C}_{k}$ is given by the fraction of the training set points assigned to that class.

Now consider the maximization with respect to $\boldsymbol{\mu}_{1}$. Again, we can pick out of the log likelihood function those terms that depend on $\boldsymbol{\mu}_{1}$:

\begin{equation}
\sum_{n=1}^{N}t_{n}\ln\mathcal{N}(\mathbf{x}_{n}|\boldsymbol{\mu}_{1},\boldsymbol{\Sigma})=-\frac{1}{2}\sum_{n=1}^{N}t_{n}(\mathbf{x}_{n}-\boldsymbol{\mu}_{1})^{\mathrm{T}}\boldsymbol{\Sigma}^{-1}(\mathbf{x}_{n}-\boldsymbol{\mu}_{1})+\text{const.}
\tag{5.57}
\end{equation}

Setting the derivative with respect to $\boldsymbol{\mu}_{1}$ to zero and rearranging, we obtain

\begin{equation}
\boldsymbol{\mu}_{1}=\frac{1}{N_{1}}\sum_{n=1}^{N}t_{n}\mathbf{x}_{n},
\tag{5.58}
\end{equation}

which is simply the mean of all the input vectors $\mathbf{x}_{n}$ assigned to class $\mathcal{C}_{1}$. By a similar argument, the corresponding result for $\boldsymbol{\mu}_{2}$ is given by

\begin{equation}
\boldsymbol{\mu}_{2}=\frac{1}{N_{2}}\sum_{n=1}^{N}(1-t_{n})\mathbf{x}_{n},
\tag{5.59}
\end{equation}

which again is the mean of all the input vectors $\mathbf{x}_{n}$ assigned to class $\mathcal{C}_{2}$.

Finally, consider the maximum likelihood solution for the shared covariance matrix $\boldsymbol{\Sigma}$. Picking out the terms in the log likelihood function that depend on $\boldsymbol{\Sigma}$, we have

\begin{align}
&-\frac{1}{2}\sum_{n=1}^{N}t_{n}\ln|\boldsymbol{\Sigma}|-\frac{1}{2}\sum_{n=1}^{N}t_{n}(\mathbf{x}_{n}-\boldsymbol{\mu}_{1})^{\mathrm{T}}\boldsymbol{\Sigma}^{-1}(\mathbf{x}_{n}-\boldsymbol{\mu}_{1}) \\
&-\frac{1}{2}\sum_{n=1}^{N}(1-t_{n})\ln|\boldsymbol{\Sigma}|-\frac{1}{2}\sum_{n=1}^{N}(1-t_{n})(\mathbf{x}_{n}-\boldsymbol{\mu}_{2})^{\mathrm{T}}\boldsymbol{\Sigma}^{-1}(\mathbf{x}_{n}-\boldsymbol{\mu}_{2}) \\
&=-\frac{N}{2}\ln|\boldsymbol{\Sigma}|-\frac{N}{2}\operatorname{Tr}\left\{\boldsymbol{\Sigma}^{-1}\mathbf{S}\right\}
\tag{5.60}
\end{align}

where we have defined

\begin{align}
\mathbf{S} &= \frac{N_{1}}{N}\mathbf{S}_{1}+\frac{N_{2}}{N}\mathbf{S}_{2} \tag{5.61}\\ 
\mathbf{S}_{1} &= \frac{1}{N_{1}}\sum_{n\in\mathcal{C}_{1}}(\mathbf{x}_{n}-\boldsymbol{\mu}_{1})(\mathbf{x}_{n}-\boldsymbol{\mu}_{1})^{\mathrm{T}} \tag{5.62}\\ 
\mathbf{S}_{2} &= \frac{1}{N_{2}}\sum_{n\in\mathcal{C}_{2}}(\mathbf{x}_{n}-\boldsymbol{\mu}_{2})(\mathbf{x}_{n}-\boldsymbol{\mu}_{2})^{\mathrm{T}}. \tag{5.63}
\end{align}

Using the standard result for the maximum likelihood solution for a Gaussian distribution, we see that $\boldsymbol{\Sigma}=\mathbf{S}$, which represents a weighted average of the covariance matrices associated with each of the two classes separately.

This result is easily extended to the $K$-class problem to obtain the corresponding maximum likelihood solutions for the parameters in which each class-conditional density is Gaussian with a shared covariance matrix. Note that the approach of fitting Gaussian distributions to the classes is not robust to outliers, because the maximum likelihood estimation of a Gaussian is not robust.

\subsection{Discrete Features}

Let us now consider discrete feature values $x_{i}$. For simplicity, we begin by looking at binary feature values $x_{i}\in\{0,1\}$ and discuss the extension to more general discrete features shortly. If there are $D$ inputs, then a general distribution would correspond to a table of $2^{D}$ numbers for each class and have $2^{D}-1$ independent variables (due to the summation constraint). Because this grows exponentially with the number of features, we can seek a more restricted representation. Here we will make the naive Bayes assumption in which the feature values are treated as independent and conditioned on the class $\mathcal{C}_{k}$. Thus, we have class-conditional distributions of the form

\begin{equation}
p(\mathbf{x}|\mathcal{C}_{k})=\prod_{i=1}^{D}\mu_{ki}^{x_{i}}(1-\mu_{ki})^{1-x_{i}},
\tag{5.64}
\end{equation}

which contain $D$ independent parameters for each class. Substituting into (5.46) then gives

\begin{equation}
a_{k}(\mathbf{x})=\sum_{i=1}^{D}\left\{x_{i}\ln\mu_{ki}+(1-x_{i})\ln(1-\mu_{ki})\right\}+\ln p(\mathcal{C}_{k}),
\tag{5.65}
\end{equation}

which again are linear functions of the input values $x_{i}$. For $K=2$ classes, we can alternatively consider the logistic sigmoid formulation given by (5.40). Analogous results are obtained for discrete variables that take $L>2$ states.

\subsection{Exponential Family}

As we have seen, for both Gaussian distributed and discrete inputs, the posterior class probabilities are given by generalized linear models with logistic sigmoid ($K=2$ classes) or softmax ($K\geqslant 2$ classes) activation functions. These are particular cases of a more general result obtained by assuming that the class-conditional densities $p(\mathbf{x}|\mathcal{C}_{k})$ are members of the subset of the exponential family of distributions given by

\begin{equation}
p(\mathbf{x}|\boldsymbol{\lambda}_{k},s)=\frac{1}{s}h\left(\frac{1}{s}\mathbf{x}\right)g(\boldsymbol{\lambda}_{k})\exp\left\{\frac{1}{s}\boldsymbol{\lambda}_{k}^{\mathrm{T}}\mathbf{x}\right\}.
\tag{5.66}
\end{equation}

Here the scaling parameter $s$ is shared across all the classes.

For the two-class problem, we substitute this expression for the class-conditional densities into (5.41) and we see that the posterior class probability is again given by a logistic sigmoid acting on a linear function $a(\mathbf{x})$, which is given by

\begin{equation}
a(\mathbf{x})=(\boldsymbol{\lambda}_{1}-\boldsymbol{\lambda}_{2})^{\mathrm{T}}\mathbf{x}+\ln g(\boldsymbol{\lambda}_{1})-\ln g(\boldsymbol{\lambda}_{2})+\ln p(\mathcal{C}_{1})-\ln p(\mathcal{C}_{2}).
\tag{5.67}
\end{equation}

Similarly, for the $K$-class problem, we substitute the class-conditional density expression into (5.46) to give

\begin{equation}
a_{k}(\mathbf{x})=\boldsymbol{\lambda}_{k}^{\mathrm{T}}\mathbf{x}+\ln g(\boldsymbol{\lambda}_{k})+\ln p(\mathcal{C}_{k})
\tag{5.68}
\end{equation}

and so again is a linear function of $\mathbf{x}$.

\end{document}
